\chapter{Performance and scale}

As repositories grow, workflows become slower and more expensive.
Performance work is about removing unnecessary work and controlling concurrency.

\section{Cancel redundant runs}
Use concurrency to cancel obsolete runs when new commits arrive.

\begin{lstlisting}[language=YAML]
concurrency:
  group: ${{ github.workflow }}-${{ github.ref }}
  cancel-in-progress: true
\end{lstlisting}

\section{Reduce the work}
\begin{itemize}[nosep]
  \item Use path filters to skip workflows when nothing relevant changed.
  \item Split expensive tests into shards.
  \item Cache dependencies aggressively.
\end{itemize}

\section{Parallelize safely}
Matrix jobs increase parallelism but also cost.
Limit with \code{max-parallel} and avoid large cartesian products.

\section{Self-hosted runners}
Self-hosted runners can be faster and cheaper for heavy workloads.
They also require patching, monitoring, and security hardening.

\section{Workflow reuse}
Reuse workflows instead of duplicating logic across repos.
Centralized workflows are easier to optimize.

\section{Measuring performance}
Track total pipeline time, queue time, and failure rate.
Optimize the slowest 20 percent of jobs first.

\section{Queue time}
Long queue times often indicate limited runners or oversized matrices.
Reduce the matrix size or add runners before optimizing job steps.

\section{Cost controls}
Use path filters, job timeouts, and max-parallel limits to control costs.
Review which workflows run on every push and reduce unnecessary triggers.

\begin{checklistbox}{Checklist: performance}
\begin{itemize}[nosep]
  \item Cancel redundant runs.
  \item Filter by paths and branches.
  \item Cache what is expensive to rebuild.
  \item Limit matrix size.
  \item Measure and prioritize the slowest jobs.
\end{itemize}
\end{checklistbox}
